\chapter{醇：活化 OH 基团}\label{ch:b1:alcohol-activation}

醇无处不在：可以大批量地由烯烃与发酵制得，
又可由上一章中的每次格氏加成得到。
然而，醇不像卤代烷那样发生取代反应：
若要取代，OH 基团必须以氢氧根离子的形式离去，而氢氧根是\omterm{def:b1:predominant-reaction:strong-acid}{强
碱}，且是极差的\omterm{def:b1:electronic-effects:nucleophile}{离去基团}（见\cref{ch:b1:nucleophilic-substitution}）。
化学家有三种办法。其一是去掉 OH 上的质子，使
\omterm{def:b1:alcohol-activation:alkoxide}{醇盐}成为强\omterm{def:b1:electronic-effects:nucleophile}{亲核试剂}；其二是加上一个质子，使\omterm{def:b1:electronic-effects:nucleophile}{离去基团}
变为水；其三是用一个基团取代氢，使氧成为
优良\omterm{def:b1:electronic-effects:nucleophile}{离去基团}的一部分，同时不触及
碳原子。本章依次讨论这三条途径，以及随之而来的
\omterm{def:b1:elementary-steps:elementary-step}{反应机理}与立体化学选择。

\begin{recall}
SN1、SN2 反应及其立体化学（见\cref{ch:b1:nucleophilic-substitution}）；
E1、E2 与\omterm{prop:b1:elimination:zaitsev}{札依采夫规则}（见\cref{ch:b1:elimination}）；
\omterm{def:b1:predominant-reaction:acidity-constant}{酸度常数}与主反应方法（见\cref{ch:b1:predominant-reaction}）；
烷基对醇酸性的影响
（见\cref{exo:b1:electronic-effects:11}）。
\end{recall}

\section{醇作为酸：醇盐}

\begin{definition}[醇盐]\label{def:b1:alcohol-activation:alkoxide}
\emph{醇盐}\index{醇盐}是醇的共轭碱 \ce{RO-}，
即 \ce{ROH} 去质子后所得，例如甲氧根 \ce{CH3O-}、乙氧根 \ce{CH3CH2O-}、
叔丁氧根 \ce{(CH3)3CO-}。醇盐是\omterm{def:b1:predominant-reaction:strong-acid}{强碱}，也是
良好的\omterm{def:b1:electronic-effects:nucleophile}{亲核试剂}。
\end{definition}

\begin{proposition}[醇的酸性]\label{prop:b1:alcohol-activation:alcohol-acidity}
醇的酸性弱于水或与水相当：$\mathrm{p}K_a$
15.5（甲醇）、15.9（乙醇）、17.1（丙-2-醇）、19.2
（2-甲基丙-2-醇），而水对应的共轭酸碱对 \ce{H2O}/\ce{OH-} 为 14.0。
因此氢氧根只能部分地使其去质子；要使其完全转化为
\omterm{def:b1:alcohol-activation:alkoxide}{醇盐}，需要更强的碱或碱金属。
\end{proposition}
% ledger: pka:methanol, pka:ethanol, pka:2-propanol, pka:tBuOH

\begin{proof}
该反应 \ce{ROH + OH- <=> RO- + H2O} 的平衡常数为 $K = 10^{14.0 -
\mathrm{p}K_a(\ce{ROH})}$，介于 $10^{-1.5}$ 与 $10^{-5}$ 之间，反应不利。
金属钠还原质子，\ce{2ROH + 2Na -> 2RO- + 2Na+ + H2}，而
氢化钠中的氢负离子则不可逆地夺取质子，
\ce{ROH + H- -> RO- + H2}：氢气逸出，反应
完全进行，与平衡常数的大小无关。
\end{proof}

\begin{inthelab}[氢化钠]
氢化钠以分散于矿物油中的灰色粉末形式出售，矿物油
可使其隔绝空气中的水分。它与水剧烈反应，
放出易燃的氢气；称量须迅速，分次
在惰性气体保护下加入干燥的醇中，氢气的冒泡可显示
去质子化的进程。过量的氢化物须在最后
缓慢加入醇加以破坏，切勿加水。
\end{inthelab}

\section{威廉森醚合成}

\begin{definition}[威廉森醚合成]\label{def:b1:alcohol-activation:williamson}
\emph{威廉森醚合成}\index{威廉森醚合成}是
制备醚 $\mathrm{R{-}O{-}R'}$，即\omterm{def:b1:alcohol-activation:alkoxide}{醇盐} \ce{RO-}
与卤代烷（或\omterm{def:b1:alcohol-activation:sulfonate}{磺酸酯}）$\mathrm{R'{-}X}$ 发生 SN2 反应。
\end{definition}

\begin{omfigure}
\setchemfig{atom sep=2.1em}
\schemestart
\chemfig{CH_3CH_2-@{o}O^{-}}\,\chemfig{Na^{+}}
\+
\chemfig{@{c}CH_3-[@{ci}]@{i}I}
\arrow{->}
\chemfig{CH_3CH_2-O-CH_3}
\+
\chemfig{Na^{+}}\,\chemfig{I^{-}}
\schemestop

\medskip
\chemmove{\draw[->, thick, shorten <=2pt, shorten >=2pt](o) .. controls +(80:8mm) and +(100:8mm) .. (c);
          \draw[->, thick, shorten <=2pt, shorten >=2pt](ci) .. controls +(90:5mm) and +(90:5mm) .. (i);}

{\small 威廉森合成甲乙醚：乙氧根进攻
碘甲烷的碳原子，进攻方向与碘原子相反（SN2）。}
\end{omfigure}

\begin{method}[选择两个反应物]\label{met:b1:alcohol-activation:williamson-choice}
醚 $\mathrm{R{-}O{-}R'}$ 可在氧的任一侧断开。
\begin{enumerate}
  \item 写出两种可能的组合：\ce{RO-} 与 $\mathrm{R'X}$，以及
        $\mathrm{R'O^-}$ 与 \ce{RX}。
  \item 保留卤素所在碳为甲基或\omterm{def:b1:electronic-effects:carbon-class}{伯碳}的组合：SN2 需要
        无位阻的碳原子。
  \item 舍弃含叔卤代烷的组合：\omterm{def:b1:alcohol-activation:alkoxide}{醇盐}是\omterm{def:b1:predominant-reaction:strong-acid}{强碱}，
        将转而发生消除（E2）。
  \item 若为仲卤代烷，则预期得到醚与烯烃的混合物。
\end{enumerate}
\end{method}

\begin{example}[不对称醚]\label{ex:b1:alcohol-activation:williamson}
2-甲氧基-2-甲基丙烷（甲基叔丁基醚）：叔丁醇钠与碘甲烷
反应可干净地得到它；甲醇钠与叔溴代烷
2-溴-2-甲基\-丙烷反应，经 E2 得到 2-甲基丙烯。体积大的
基团放在\omterm{def:b1:alcohol-activation:alkoxide}{醇盐}一侧，小基团放在卤代烷一侧。
\end{example}

\section{活化 OH 基团}

\begin{definition}[磺酸酯]\label{def:b1:alcohol-activation:sulfonate}
\emph{磺酸酯}\index{磺酸酯} $\mathrm{R{-}O{-}SO_2{-}R'}$ 由醇与
$\mathrm{R'SO_2Cl}$ 在
碱（吡啶）存在下生成。\emph{对甲苯磺酸酯}\index{对甲苯磺酸酯}\ce{R-OTs}由
4-甲基苯磺酰氯（对甲苯磺酰氯）制得；
\emph{甲磺酸酯}\index{甲磺酸酯}\ce{R-OMs}，由甲磺酰氯制得。
磺酸根离子 $\mathrm{R'SO_3^-}$ 是\omterm{def:b1:predominant-reaction:strong-acid}{强酸}的共轭碱，
并经共振在三个氧原子上得到稳定，是优良的\omterm{def:b1:electronic-effects:nucleophile}{离去基团}。
\end{definition}

\begin{omfigure}
\setchemfig{atom sep=2em}
\schemestart
\chemfig{R-O-H}
\+
\chemfig{Cl-S(=[2]O)(=[6]O)-[:0]*6(-=-(-CH_3)=-=)}
\arrow{->[吡啶]}
\chemfig{R-O-S(=[2]O)(=[6]O)-[:0]*6(-=-(-CH_3)=-=)}
\schemestop
\par\vspace{6pt}

{\small 醇的对甲苯磺酰化。O--H 键被换成 O--S 键，醇中的
C--O 键不受影响。吡啶接收生成的 HCl。}
\end{omfigure}

\begin{proposition}[通过对甲苯磺酰化和取代确定构型]\label{prop:b1:alcohol-activation:tosylation-retention}
醇在\omterm{def:b1:stereochemistry-in-depth:stereocentre}{立体中心}碳上发生对甲苯磺酰化时，\omterm{def:b1:stereochemistry-in-depth:isomers}{构型}保持不变；
随后\omterm{def:b1:alcohol-activation:sulfonate}{对甲苯磺酸酯}的 SN2 取代使其翻转。
两步合起来，以翻转的方式将 OH 换成\omterm{def:b1:electronic-effects:nucleophile}{亲核试剂}。
\end{proposition}

\begin{proof}
在对甲苯磺酰化中，发生变化的键是 O--H 键与 S--Cl 键：醇的氧原子进攻硫原子，
其氢转移给碱。与\omterm{def:b1:stereochemistry-in-depth:stereocentre}{立体中心}碳相连的键没有断裂，
因此其周围的排列保持不变。
在 SN2 步骤中，C--O 键断裂，同时\omterm{def:b1:electronic-effects:nucleophile}{亲核试剂}从相反的一侧成键
（\cref{prop:b1:nucleophilic-substitution:sn2-inversion}）。
\end{proof}

\begin{omfigure}
\begin{tikzpicture}[font=\small, >={Stealth}]
  \node (a) at (0,0) {\chemfig{C(-[:150]H_3C)(<[:210]H)(<:[:250]C_2H_5)-OH}};
  \node at (0,-1.4) {\cip{S}-丁-2-醇};
  \draw[->, thick] (1.3,0) -- (2.5,0) node[midway, above] {对甲苯磺酰氯} node[midway, below, font=\scriptsize] {构型保持};
  \node (b) at (3.9,0) {\chemfig{C(-[:150]H_3C)(<[:210]H)(<:[:250]C_2H_5)-OTs}};
  \node at (3.9,-1.4) {\cip{S} 对甲苯磺酸酯};
  \draw[->, thick] (5.3,0) -- (6.6,0) node[midway, above] {\ce{CN-}} node[midway, below, font=\scriptsize] {SN2，翻转};
  \node (c) at (8.3,0) {\chemfig{[:0]NC-C(-[:30]CH_3)(<[:-30]H)(<:[:-70]C_2H_5)}};
  \node at (8.3,-1.4) {\cip{R} 腈};
\end{tikzpicture}

{\small 从 \cip{S}-丁-2-醇转化为 2-甲基丁腈：对甲苯磺酰化
保持\omterm{def:b1:stereochemistry-in-depth:isomers}{构型}，与氰根发生的 SN2 步骤使其翻转。总体：
OH 被 CN 以翻转的方式取代。}
\end{omfigure}

同一目标也可在酸性条件下达到：OH 基团质子化后以
水的形式离去，\ce{R-OH2+ -> R+ + H2O}（SN1，E1），或受进攻（SN2）。
然而，酸性条件只把\omterm{def:b1:electronic-effects:nucleophile}{亲核试剂}限于能耐受它的那些——卤离子、水、
醇——并打开通向\omterm{def:b1:electronic-effects:intermediates}{碳正离子}及其副反应的大门。

\section{转化为卤代烷}

\begin{proposition}[醇与卤化氢]\label{prop:b1:alcohol-activation:hx-by-class}
叔醇在室温下与浓盐酸反应生成卤代烷
（SN1）；伯醇需要氢溴酸或
氢碘酸并加热（对质子化醇的 SN2）；仲醇
的\omterm{def:b1:rate-laws:rate}{反应速率}居中，常常两种机理兼有。
\end{proposition}

\begin{proof}
质子化使 OH 变为 \ce{-OH2+}。对叔醇而言，水很容易离去，
生成稳定的\omterm{def:b1:electronic-effects:intermediates}{碳正离子}，并被卤离子捕获（SN1）：即使是较弱的\omterm{def:b1:electronic-effects:nucleophile}{亲核试剂}
\ce{Cl-} 也足够。伯\omterm{def:b1:electronic-effects:intermediates}{碳正离子}过于
不稳定，因此卤离子须从背面将水推出（SN2），这需要
较好的\omterm{def:b1:electronic-effects:nucleophile}{亲核试剂}（\ce{Br-}、\ce{I-}，在质子性介质中比 \ce{Cl-} 更好，
见\cref{prop:b1:electronic-effects:nucleophilicity-basicity}）以及加热。
\end{proof}

\begin{example}[一个分类试验]\label{ex:b1:alcohol-activation:lucas}
浓盐酸与氯化锌（一种有助于 OH 离去的\omterm{def:b1:lewis-resonance-vsepr:lewis-acid}{路易斯酸}）的混合物
使叔醇立即变浑浊——不溶性的氯代烷析出——
仲醇在数分钟内变浑浊，
而伯醇在室温下几乎不变：
\omterm{def:b1:electronic-effects:intermediates}{碳正离子}稳定性的顺序在试管中得以直观呈现。
\end{example}

\section{酸中的脱水与醚的生成}

\begin{definition}[醇的脱水]\label{def:b1:alcohol-activation:dehydration}
醇的\emph{脱水}\index{脱水}是指
消去水分子的反应，在\omterm{def:b1:predominant-reaction:strong-acid}{强酸}（硫酸或磷酸）催化并加热下进行，
生成烯烃：$\mathrm{R_2CH{-}CR'_2OH} \to \mathrm{R_2C{=}CR'_2} + \ce{H2O}$。
\end{definition}

\begin{omfigure}
\setchemfig{atom sep=2.1em}
\schemestart
\chemfig{C(-[2]CH_3)(-[:210]H_3C)(-[:-30]CH_2CH_3)-OH}
\arrow{->[\ce{H+}]}
\chemfig{C(-[2]CH_3)(-[:210]H_3C)(-[:-30]CH_2CH_3)-OH_2^{+}}
\arrow{->[\ce{-H2O}]}
\chemfig{C^{+}(-[2]CH_3)(-[:210]H_3C)(-[:-30]CH_2CH_3)}
\schemestop

\medskip
\schemestart
\chemfig{C^{+}(-[2]CH_3)(-[:210]H_3C)(-[:-30]CH_2CH_3)}
\arrow{->[\ce{-H+}][（主）]}
\chemfig{C(-[2]CH_3)(-[:210]H_3C)=[:-30]CHCH_3}
\schemestop

\medskip
{\small 2-甲基丁-2-醇的酸催化\omterm{def:b1:alcohol-activation:dehydration}{脱水}（E1）：质子化，
失去水生成叔\omterm{def:b1:electronic-effects:intermediates}{碳正离子}，再失去一个质子。从
\ce{CH2} 上失去质子得到三取代的 2-甲基丁-2-烯
（\omterm{prop:b1:elimination:zaitsev}{札依采夫规则}的产物，为主）；从甲基上失去质子得到 2-甲基丁-1-烯（次要）。}
\end{omfigure}

\begin{proposition}[醚还是烯烃]\label{prop:b1:alcohol-activation:ether-vs-alkene}
伯醇与硫酸共热时，
在中等温度下主要生成对称醚（一分子通过 SN2 取代另一分子），
在较高温度下主要生成烯烃；
仲醇和叔醇只生成烯烃。
\end{proposition}

\begin{proof}
对于伯醇，质子化的醇既可以在
碳原子上被第二个醇分子进攻（SN2，失去一个质子后生成 \ce{R-O-R}），
也可以在 $\beta$ 位氢处发生消除（E2）。
消除生成的分子多于消耗的分子，因此随温度升高而更有利，
（见\cref{prop:b1:elimination:sn-vs-e}）；温度较低时，
则以取代为主。仲醇和叔醇电离生成\omterm{def:b1:electronic-effects:intermediates}{碳正离子}，它
失去质子比被弱\omterm{def:b1:electronic-effects:nucleophile}{亲核试剂}捕获更容易，加热时更是如此；
生成的水还可借助
挥发性烯烃的蒸馏而被移去。
\end{proof}

\section{练习}

\begin{exercise}[$\star$]\label{exo:b1:alcohol-activation:1}
写出乙醇分别与钠、氢化钠
和氢氧化钠的反应；并计算后者的常数（乙醇的 $\mathrm{p}K_a$
为 15.9）。
\end{exercise}

\begin{exercise}[$\star$]\label{exo:b1:alcohol-activation:2}
写出下列反应的产物：甲醇钠与 1-溴丙烷；乙醇钠
与碘甲烷；苯酚钠与溴乙烷。
\end{exercise}

\begin{exercise}[$\star$]\label{exo:b1:alcohol-activation:3}
写出丙-1-醇的对甲苯磺酰化反应，以及\omterm{def:b1:alcohol-activation:sulfonate}{对甲苯磺酸酯}
与碘化钠的反应；这一途径相对于醇直接与碘化钠
反应有何优点？
\end{exercise}

\begin{exercise}[$\star$]\label{exo:b1:alcohol-activation:4}
写出丁-2-醇和
2-甲基戊-2-醇经酸催化\omterm{def:b1:alcohol-activation:dehydration}{脱水}所生成的烯烃，并指出其中的主要产物。
\end{exercise}

\begin{exercise}[$\star\star$]\label{exo:b1:alcohol-activation:5}
提出合成 1-乙氧基-2-甲基丙烷的\omterm{def:b1:alcohol-activation:williamson}{威廉森醚合成}方案，并说明
所选两种反应物的理由。若选另一种搭配，为何会得到烯烃？
\end{exercise}

\begin{exercise}[$\star\star$]\label{exo:b1:alcohol-activation:6}
\cip{R}-辛-2-醇转化为其\omterm{def:b1:alcohol-activation:sulfonate}{对甲苯磺酸酯}，然后
用乙醇钠的乙醇溶液处理。写出所得醚的\omterm{def:b1:stereochemistry-in-depth:isomers}{构型}。
若将同一醇与硫酸共热，将得到什么？
\end{exercise}

\begin{exercise}[$\star\star$]\label{exo:b1:alcohol-activation:7}
写出 2-甲基丙-2-醇与
浓盐酸反应的机理，并说明为何该反应
在室温下进行得很快。
\end{exercise}

\begin{exercise}[$\star\star$]\label{exo:b1:alcohol-activation:8}
写出乙醇在硫酸中生成
乙醚的机理，以及与之竞争的消除反应。
\end{exercise}

\begin{exercise}[$\star\star$]\label{exo:b1:alcohol-activation:9}
在本章的鉴别实验中，丁-1-醇、丁-2-醇和
2-甲基丙-2-醇中哪一种最先变浑浊？写出
该反应的产物。
\end{exercise}

\begin{exercise}[$\star\star\star$]\label{exo:b1:alcohol-activation:10}
3,3-二甲基丁-2-醇在酸中加热，主要得到
2,3-二甲基丁-2-烯，其碳骨架与原醇的不同。
请提出一个机理，其中\omterm{def:b1:electronic-effects:intermediates}{碳正离子}内的甲基发生
1,2-迁移，并说明这一迁移的理由。
\end{exercise}

\begin{exercise}[$\star\star\star$]\label{exo:b1:alcohol-activation:11}
设计 \cip{S}-2-甲氧基丁烷的合成路线，以 \cip{R}-丁-2-醇为原料，
另一条由 \cip{S}-丁-2-醇出发，两条路线都要
使\omterm{def:b1:stereochemistry-in-depth:isomers}{构型}得到控制。
\end{exercise}

\begin{exercise}[$\star\star\star$]\label{exo:b1:alcohol-activation:12}
证明平衡 \ce{2CH3CH2OH <=> (CH3CH2)2O + H2O} 可以
通过连续移除水而朝醚的方向进行，并利用
\omterm{def:b1:extent-q-and-k:quotient}{反应商}（见\cref{ch:b1:extent-q-and-k}）加以说明。
\end{exercise}

\section{问题：通往醚的两条路线}

\begin{problem}[{周末问题——通往甲基叔丁基醚的两条威廉森路线，
由旋光性醇制得的立体化学可控的醚，
叔醇的脱水，以及所得醚的质
量}]
\label{pb:b1:alcohol-activation:1}
2-甲氧基-2-甲基丙烷（甲基叔丁基醚，MTBE）曾
大规模生产，用作燃料添加剂。摩尔质量（\unit{g/mol}）：\ce{C4H10O}
74.0，\ce{C5H12O} 88.0，\ce{CH3I} 141.9。

\medskip
\noindent\textbf{第一部分——用威廉森合成法制备 MTBE。}
\begin{enumerate}
  \item 写出可生成 MTBE 的两对反应物
        （\omterm{def:b1:alcohol-activation:alkoxide}{醇盐}与卤代烷）。
  \item 哪一对反应物不能得到 MTBE？写出实际发生的反应。
  \item 写出能够成功进行的反应及其机理。
  \item 如何由 2-甲基丙-2-醇制备叔丁醇钠？
        为何不用氢氧化钠？
  \item 为何该反应要在无水醇或\omterm{def:b1:intermolecular-forces-solvents:solvent-classes}{非质子性溶剂}中进行？
  \item 产物是否具有\omterm{def:b1:stereochemistry-in-depth:stereocentre}{手性}？
\end{enumerate}

\medskip
\noindent\textbf{第二部分——具有\omterm{def:b1:stereochemistry-in-depth:optical-activity}{旋光性}的醚。}
\begin{enumerate}[resume]
  \item \cip{R}-丁-2-醇转化为其\omterm{def:b1:alcohol-activation:sulfonate}{对甲苯磺酸酯}，写出
        该反应，并给出\omterm{def:b1:alcohol-activation:sulfonate}{对甲苯磺酸酯}的\omterm{def:b1:stereochemistry-in-depth:isomers}{构型}。
  \item 该\omterm{def:b1:alcohol-activation:sulfonate}{对甲苯磺酸酯}与甲醇钠反应，写出产物及其
        \omterm{def:b1:stereochemistry-in-depth:isomers}{构型}。
  \item 醇的哪一种\omterm{def:b1:stereochemistry-in-depth:isomers}{构型}经同一路线可以得到
        \cip{R}-2-甲氧基丁烷？
  \item 另一条路线：\cip{R}-丁-2-醇的\omterm{def:b1:alcohol-activation:alkoxide}{醇盐}与
        碘甲烷反应，产物的\omterm{def:b1:stereochemistry-in-depth:isomers}{构型}如何？为什么？
  \item 第 8 小题中与之竞争的副反应是什么，
        以及它为何在此处受到限制？
  \item 若将 \cip{R}-丁-2-醇与甲醇在硫酸中加热，
        为何会得到接近外消旋的醚（若有的话）？
\end{enumerate}

\medskip
\noindent\textbf{第三部分——叔醇的\omterm{def:b1:alcohol-activation:dehydration}{脱水}。}
\begin{enumerate}[resume]
  \item 写出 2-甲基丁-2-醇
        酸催化\omterm{def:b1:alcohol-activation:dehydration}{脱水}的机理。
  \item 写出两种烯烃，并预测主要产物。
  \item 为何该醇不生成醚？
  \item 为何 2-甲基丁-2-烯一经生成便被蒸出？
  \item 画出 E1 步骤的\omterm{def:b1:elementary-steps:energy-profile}{能量曲线}，并标出
        \omterm{def:b1:elementary-steps:rds}{决速步}。
  \item 伯醇能否经同样的机理\omterm{def:b1:alcohol-activation:dehydration}{脱水}？
\end{enumerate}

\medskip
\noindent\textbf{第四部分——质量计算。}
\begin{enumerate}[resume]
  \item 将 \qty{10.0}{g} 的 2-甲基丙-2-醇转化为
        \omterm{def:b1:alcohol-activation:alkoxide}{醇盐}，并计算其物质的量。
  \item 要使碘甲烷过量 \qty{10}{\percent}，需要多少质量的碘甲烷？
  \item 计算 MTBE 的理论质量。
  \item 计算\omterm{def:b1:extent-q-and-k:fractional-extent}{产率}为 \qty{75}{\percent} 时所得的质量。
\end{enumerate}
\end{problem}
